\section{LITERATURE REVIEW}\label{SC:LR}\label{sec:literature}
\begin{revised}
The literature is grouped below by the information used for learning and
the role of verification in each approach.
Early statistical and data-mining reviews distinguish label-based detection
from anomaly screening \citep{Bolton2002,Ngai2011}; broader financial-fraud
reviews describe their continuing complementarity \citep{West2016,A61}.
The systematic review of \citet{Schrijver2024} describes the prominence of
supervised learning and oversampling in automobile-insurance fraud studies,
as well as limitations in public data availability.

\subsection{Supervised approaches}
Supervised models use previously investigated claims to estimate a decision
boundary or ranking. Regression, support-vector machines, forests, boosting,
and neural networks are common choices \citep{A80,A61}. Insurance-specific
work includes optimized boosting and interpretable assessment
\citep{Ding2025}, practical model-selection guidance
\citep{BanulescuRadu2024}, and imbalance and missing-value handling
\citep{Khalil2024}. These methods can exploit claim attributes effectively
when their training labels are sufficiently representative. Direct comparisons
of supervised and unsupervised insurance models are also reported by
\citet{Debener2023}. Selective
historical investigation, however, can bias which examples are labeled.

In \emph{Annals of Data Science}, \citet{Kamalov2023} examine feature selection
under imbalance and propose a minority-sensitive criterion. Their work
motivates explicit attention to fraud-class metrics here; we do not equate
our unsupervised Gini screen with their supervised selection algorithm.

\subsection{Unsupervised approaches}
Anomaly methods rank observations according to deviations from a fitted
notion of normality. Auto-insurance applications include spectral ranking
\citep{Nian2016} and unsupervised deep learning \citep{Gomes2021}.
Such approaches can propose cases before reliable labels exist, but unusual
observations need not be fraudulent. Clustering methods also depend on
representation and dissimilarity choices, as reviewed by
\citet{XuTian2015} in \emph{Annals of Data Science}. This supports using
heterogeneous signals while retaining verification of their interpretation.

\subsection{Reinforcement-learning-based approaches}
When verification is costly, a learned policy can decide which observations
justify acquiring an outcome. In the framework studied by \citet{A73},
reinforcement learning guides this information-acquisition decision for
outlier detection. Actions are evaluated through a reward signal. SAD instead uses fixed voting
thresholds and a random audit fraction; it has no learned reward-maximizing
query policy. Its sequential component is retraining after each reviewed
batch, which can be tested against a frozen model.

\subsection{Hybrid architectures}
Hybrid approaches combine representations or learners, while graph methods
use relationships among claims, people, or entities \citep{Vorobyev2024}.
PSO-XGBoost is a recent insurance-specific example of optimization combined
with boosting \citep{Ding2025}. \citet{Hassani2020}, in \emph{Annals of Data
Science}, review deep-learning applications in banking, providing broader
financial-domain context. That review is contextual evidence, not a direct
insurance-fraud performance comparator.

\subsection{Research gap and comparator scope}
SAD addresses the connection between candidate generation, review, and
iterative learning. It does not replace the need for representative labels
or resolve all challenges in recent deep or graph-based systems.
Table \ref{tab:positioning} distinguishes these roles. The numerical
experiments include an explicitly adapted PSO-XGBoost comparator and balanced
random forests in addition to eight component classifiers. We do not
attribute our adapted implementation's numbers to the original publication.

\begin{table}[tbp]\centering\footnotesize
\caption{Comparative positioning by learning and acquisition mechanism.}
\label{tab:positioning}
\begin{tabularx}{\linewidth}{>{\raggedright\arraybackslash}p{2.5cm}*{3}{>{\raggedright\arraybackslash}X}}
\toprule
Family & Representative source & Strength & Boundary relevant to SAD\\\midrule
Supervised & Ding et al.; Banulescu-Radu and Yankol-Schalck & Direct label-based discrimination & Needs representative verified labels\\
Unsupervised & Nian et al.; Gomes et al. & Initial screening without outcomes & Anomaly is not confirmed fraud\\
Reinforcement learning & Nijhuis and van Lelyveld & Values costly information & Requires a query/reward design\\
Hybrid/graph & Vorobyev; Ding et al. & Relational or optimized predictive signals & May require links or additional tuning\\
SAD & This study & Explicit review and batch update workflow & Retrospective simulation; policy-dependent labels\\
\bottomrule
\end{tabularx}
\end{table}

\begin{table}[tbp]\centering\footnotesize
\caption{Mechanisms of recent counterparts compared with SAD.
Design comparison does not imply numerical reproduction.}
\label{tab:mechanisms}
\begin{tabularx}{\linewidth}{>{\raggedright\arraybackslash}p{2.8cm}*{3}{>{\raggedright\arraybackslash}X}}\toprule
Method & Main mechanism & Relation to sequential review & Relation to multilayer SAD\\\midrule
PSO-XGBoost \citep{Ding2025} & Optimized boosted trees & Search tunes parameters; it does not implement batch review & No post-consensus entropy gate specified\\
Unsupervised deep learning \citep{Gomes2021} & Learned anomaly representation & Does not by itself define SAD's verified-label accumulation & Different initial scoring mechanism\\
Costly-verification RL \citep{A73} & Reward-driven information acquisition & Learns an acquisition policy; SAD uses fixed rules & SAD's entropy gate is not a reward policy\\
Claims graph features \citep{Vorobyev2024} & Relational feature construction & Requires links beyond independent claim attributes & Adds relational information rather than SAD's entropy gate\\
SAD, this study & Anomaly votes, verification, supervised votes & Retrains once after each reviewed batch & Applies candidate-local EAD after initial consensus\\
\bottomrule\end{tabularx}\end{table}
\end{revised}
